\documentclass{article}
\usepackage{amsmath,amssymb,amsthm}
\usepackage[ruled,vlined]{algorithm2e}
\usepackage{enumitem}
\usepackage{hyperref}
\newtheorem{theorem}{Theorem}
\newtheorem{lemma}{Lemma}
\newtheorem{assumption}{Assumption}
\newtheorem{corollary}{Corollary}
\newtheorem{definition}{Definition}
\newcommand{\R}{\mathbb{R}}
\newcommand{\Cset}{\mathcal{C}}
\newcommand{\grad}{\nabla}
\newcommand{\argmin}{\operatorname*{arg\,min}}
\newcommand{\argmax}{\operatorname*{arg\,max}}
\newcommand{\inner}[1]{\langle #1\rangle}
\begin{document}

\section*{Algorithm Name}
\textbf{Away‑Step Frank–Wolfe (AFW)} for smooth convex minimization over a polytope.

\section*{Mathematical Formulation}
Let $f:\R^{d}\rightarrow\R$ be convex and $L$‑smooth.  
Let $\Cset\subset\R^{d}$ be a compact convex polytope with vertex set $\mathcal{V}=\{v^{1},\dots ,v^{m}\}$.  
The problem is
\[
\min_{x\in\Cset} f(x). \tag{P}
\]

Given an iterate $x_{k}\in\Cset$ we maintain a \emph{representation}
\[
x_{k}= \sum_{v\in\mathcal{S}_{k}} \alpha_{k}^{(v)} v ,\qquad
\alpha_{k}^{(v)}\ge 0,\ \sum_{v\in\mathcal{S}_{k}}\alpha_{k}^{(v)}=1,
\]
where $\mathcal{S}_{k}\subseteq\mathcal{V}$ is the \emph{active set}.

\begin{enumerate}[label=\textbf{Step \arabic*:}, leftmargin=*]
\item \textbf{FW direction.} Compute the linear minimization oracle (LMO)
\[
s_{k}\;:=\;\argmin_{s\in\Cset}\;\inner{\grad f(x_{k})}{s}.
\]
The Frank–Wolfe direction is $d_{k}^{\text{FW}}:=s_{k}-x_{k}$.

\item \textbf{Away direction.} Choose the vertex in the active set that maximizes the linear
approximation:
\[
v_{k}\;:=\;\argmax_{v\in\mathcal{S}_{k}}\;\inner{\grad f(x_{k})}{v},
\qquad
d_{k}^{\text{A}}:=x_{k}-v_{k}.
\]

\item \textbf{Direction selection.} Compute the directional derivatives
\[
\Delta^{\text{FW}}_{k}:=\inner{\grad f(x_{k})}{d_{k}^{\text{FW}}},\qquad
\Delta^{\text{A}}_{k}:=\inner{\grad f(x_{k})}{d_{k}^{\text{A}}}.
\]
If $\Delta^{\text{FW}}_{k}\le \Delta^{\text{A}}_{k}$ set $d_{k}=d_{k}^{\text{FW}}$ (a \emph{FW step});
otherwise set $d_{k}=d_{k}^{\text{A}}$ (an \emph{away step}).

\item \textbf{Step‑size.} Let
\[
\gamma_{\max }:=\begin{cases}
1, & d_{k}=d_{k}^{\text{FW}},\\[4pt]
\displaystyle\frac{\alpha_{k}^{(v_{k})}}{1-\alpha_{k}^{(v_{k})}}, & d_{k}=d_{k}^{\text{A}} .
\end{cases}
\]
Choose $\gamma_{k}\in[0,\gamma_{\max}]$ by exact line search:
\[
\gamma_{k}\;:=\;\argmin_{\gamma\in[0,\gamma_{\max}]}\; f\!\bigl(x_{k}+\gamma d_{k}\bigr).
\]

\item \textbf{Update.}
\[
x_{k+1}=x_{k}+\gamma_{k}d_{k}.
\]
If a FW step was taken, add $s_{k}$ to the active set (or increase its coefficient);
if an away step was taken, decrease the coefficient of $v_{k}$ and, when it becomes
zero, remove $v_{k}$ from $\mathcal{S}_{k}$.
\end{enumerate}

\section*{Pseudocode}
\begin{algorithm}[H]
\DontPrintSemicolon
\SetAlgoLined
\KwIn{Convex $L$‑smooth $f$, polytope $\Cset=\operatorname{conv}(\mathcal{V})$, tolerance $\varepsilon>0$}
\KwOut{Approximate solution $x$}
Choose any vertex $v^{(i)}\in\mathcal{V}$ and set $x_{0}=v^{(i)}$, $\mathcal{S}_{0}=\{v^{(i)}\}$,
$\alpha^{(v^{(i)})}_{0}=1$\;
\For{$k=0,1,2,\dots$}{
    $\displaystyle g_{k}\gets\grad f(x_{k})$\;
    $s_{k}\gets\argmin_{s\in\Cset}\inner{g_{k}}{s}$\;
    $v_{k}\gets\argmax_{v\in\mathcal{S}_{k}}\inner{g_{k}}{v}$\;
    $d^{\text{FW}}_{k}\gets s_{k}-x_{k}$\;
    $d^{\text{A}}_{k}\gets x_{k}-v_{k}$\;
    $\Delta^{\text{FW}}_{k}\gets\inner{g_{k}}{d^{\text{FW}}_{k}}$,\;
    $\Delta^{\text{A}}_{k}\gets\inner{g_{k}}{d^{\text{A}}_{k}}$\;
    \eIf{$\Delta^{\text{FW}}_{k}\le\Delta^{\text{A}}_{k}$}{
        $d_{k}\gets d^{\text{FW}}_{k}$,\;
        $\gamma_{\max}\gets1$\;
    }{
        $d_{k}\gets d^{\text{A}}_{k}$,\;
        $\gamma_{\max}\gets\displaystyle\frac{\alpha^{(v_{k})}_{k}}{1-\alpha^{(v_{k})}_{k}}$\;
    }
    $\displaystyle\gamma_{k}\gets\argmin_{\gamma\in[0,\gamma_{\max}]} f(x_{k}+\gamma d_{k})$\;
    $x_{k+1}\gets x_{k}+\gamma_{k}d_{k}$\;
    Update the coefficients $\{\alpha^{(\cdot)}_{k}\}$ and the active set $\mathcal{S}_{k}$ accordingly\;
    \If{$-\Delta^{\text{FW}}_{k}\le\varepsilon$}{break}
}
\Return{$x_{k+1}$}
\caption{Away‑Step Frank–Wolfe (AFW)}
\end{algorithm}

\section*{Dry Run Example}
Minimize $f(w)=(w-3)^{2}$ over the 1‑dimensional polytope $\Cset=[0,5]$.
The vertex set is $\mathcal{V}=\{0,5\}$.
We start at $w_{0}=0$ (active set $\mathcal{S}_{0}=\{0\}$).

\begin{center}
\begin{tabular}{c|c|c|c|c}
Iter & $g_{k}=\grad f(w_{k})$ & $s_{k}$ (LMO) & $\gamma_{k}$ (line search) & $w_{k+1}$ \\ \hline
0 & $-6$ & $\argmin_{s\in\{0,5\}}(-6)s=5$ & $\displaystyle\argmin_{\gamma\in[0,1]}(5\gamma-3)^{2}=0.6$ & $0+0.6\cdot5=3$ \\ \hline
1 & $0$ & any vertex (choose $5$) & $\gamma_{1}=0$ (already optimal) & $3$ \\ \hline
2 & $0$ & – & – & $3$ (termination) \\
\end{tabular}
\end{center}
After the first iteration the algorithm reaches the optimum $w^{\star}=3$ and
terminates.

\section*{Assumptions}
\begin{assumption}[Smoothness]\label{ass:smooth}
$f$ is $L$‑smooth on $\Cset$, i.e.
\[
\|\grad f(x)-\grad f(y)\| \le L\|x-y\|,\qquad\forall x,y\in\Cset .
\]
\end{assumption}
\begin{assumption}[Strong Convexity]\label{ass:strong}
$f$ is $\mu$‑strongly convex on $\Cset$:
\[
f(y) \ge f(x)+\inner{\grad f(x)}{y-x}+\frac{\mu}{2}\|y-x\|^{2},
\qquad\forall x,y\in\Cset .
\]
\end{assumption}
\begin{assumption}[Polytope Geometry]\label{ass:polytope}
$\Cset=\operatorname{conv}(\mathcal{V})$ is a compact polytope with
\begin{itemize}
\item diameter $D:=\max_{u,v\in\Cset}\|u-v\|$,
\item pyramidal width $\delta>0$ (see~\cite{lacoste2015linear}).
\end{itemize}
\end{assumption}
All constants $L,\mu,D,\delta$ are known to be finite and positive.

\section*{Main Convergence Theorem}
\begin{theorem}[Linear convergence of AFW]\label{thm:linear}
Under Assumptions~\ref{ass:smooth}--\ref{ass:polytope}, the iterates $\{x_{k}\}$ generated
by the Away‑Step Frank–Wolfe algorithm satisfy
\[
f(x_{k})-f(x^{\star})\;\le\;\bigl(1-\rho\bigr)^{k}\,\bigl(f(x_{0})-f(x^{\star})\bigr),
\qquad\text{with}\qquad
\rho:=\frac{\mu\,\delta^{2}}{L\,D^{2}} \in (0,1).
\tag{1}
\]
Consequently $x_{k}\to x^{\star}$ linearly.
\end{theorem}

\section*{Theoretical Properties}
\begin{itemize}
\item \textbf{Stability.} The algorithm never leaves $\Cset$ and the coefficients
remain in the simplex, guaranteeing numerical stability.
\item \textbf{Hyper‑parameter sensitivity.} AFW does not require any step‑size
tuning; the line‑search is exact and $\gamma_{\max}$ depends only on the current
active‑set coefficients.
\item \textbf{Comparison.} Classical Frank–Wolfe obtains $O(1/k)$ sublinear
rates on the same class of problems, whereas AFW achieves the linear rate
\eqref{1} thanks to away steps that eliminate ``bad'' vertices.
\end{itemize}

\section*{Discussion}
The linear rate hinges on the geometric constant $\delta$ (pyramidal width);
for simplices $\delta=1$ and the bound reduces to the classical
$\rho=\mu/(L D^{2})$.  When $f$ is only convex (no strong convexity) the
algorithm reverts to the $O(1/k)$ guarantee of the vanilla Frank–Wolfe method.

\section*{Preliminaries (Definitions and Lemmas)}
\begin{definition}[Curvature constant]\label{def:curv}
For an $L$‑smooth $f$ on $\Cset$ define
\[
C_{f}:=\sup_{\substack{x,s\in\Cset\\ \gamma\in[0,1]}}
\frac{2}{\gamma^{2}}\bigl(f(x+\gamma(s-x))-f(x)-\gamma\inner{\grad f(x)}{s-x}\bigr).
\]
For $L$‑smooth functions one has $C_{f}\le L D^{2}$.
\end{definition}
\begin{lemma}[Descent lemma]\label{lem:descent}
For any $x\in\Cset$, direction $d$ with $x+\gamma d\in\Cset$ and $\gamma\in[0,1]$,
\[
f(x+\gamma d)\le f(x)+\gamma\inner{\grad f(x)}{d}+\frac{L}{2}\gamma^{2}\|d\|^{2}.
\tag{2}
\]
\end{lemma}
\begin{proof}
Directly from $L$‑smoothness (integrate the gradient Lipschitz bound). \hfill $\square$
\end{proof}
\begin{lemma}[Geometric lower bound]\label{lem:geom}
For any $x\in\Cset$ and optimal $x^{\star}\in\Cset$,
\[
\inner{\grad f(x)}{x-x^{\star}} \;\ge\; \frac{\mu\,\delta^{2}}{D^{2}}\,
\bigl(f(x)-f(x^{\star})\bigr).
\tag{3}
\]
\end{lemma}
\begin{proof}
By strong convexity,
$f(x)-f(x^{\star})\le\inner{\grad f(x)}{x-x^{\star}}-\frac{\mu}{2}\|x-x^{\star}\|^{2}$.
Rearranging gives
$\inner{\grad f(x)}{x-x^{\star}}\ge
f(x)-f(x^{\star})+\frac{\mu}{2}\|x-x^{\star}\|^{2}$.
The definition of pyramidal width $\delta$ yields
$\|x-x^{\star}\|\ge\frac{\delta}{D}\|x-x^{\star}\|_{1}$,
and because $x$ is a convex combination of vertices,
$\|x-x^{\star}\|_{1}\ge 2\bigl(f(x)-f(x^{\star})\bigr)/L$ (using curvature bound).
Combining the inequalities leads to \eqref{eq:3}.  ∎
\end{proof}

\section*{Main Proof (Step‑by‑Step Derivation)}
We prove Theorem~\ref{thm:linear}.  Throughout the proof let $g_{k}:=\grad f(x_{k})$.
Denote the chosen direction by $d_{k}$ and the step size by $\gamma_{k}$.

\noindent\textbf{[Step 1]} \emph{Descent after the line search}.  
From Lemma~\ref{lem:descent} with $\|d_{k}\|\le D$ and using the exact line search,
\[
f(x_{k+1})\le f(x_{k})+\gamma_{k}\inner{g_{k}}{d_{k}}+\frac{L}{2}\gamma_{k}^{2}\|d_{k}\|^{2}.
\tag{4}
\]

\noindent\textbf{[Step 2]} \emph{Upper bound on the optimal step size}.  
Because $\gamma_{k}$ minimizes the right‑hand side of \eqref{4} over
$[0,\gamma_{\max}]$, the derivative at $\gamma_{k}$ (when $\gamma_{k}\in(0,\gamma_{\max})$) is zero:
\[
\inner{g_{k}}{d_{k}}+L\gamma_{k}\|d_{k}\|^{2}=0
\quad\Longrightarrow\quad
\gamma_{k}= -\frac{\inner{g_{k}}{d_{k}}}{L\|d_{k}\|^{2}}.
\tag{5}
\]
If the minimizer lies at the boundary, $\gamma_{k}=\gamma_{\max}$, inequality \eqref{5}
still holds with “$\le$’’ instead of equality.

\noindent\textbf{[Step 3]} \emph{Plugging \eqref{5} into \eqref{4}}.  
Using \eqref{5} we obtain
\[
\begin{aligned}
f(x_{k+1})
&\le f(x_{k})-\frac{\inner{g_{k}}{d_{k}}^{2}}{2L\|d_{k}\|^{2}}  \\
&\le f(x_{k})-\frac{1}{2L D^{2}}\,\inner{g_{k}}{d_{k}}^{2},
\end{aligned}
\tag{6}
\]
where we used $\|d_{k}\|\le D$.

\noindent\textbf{[Step 4]} \emph{Relating $\inner{g_{k}}{d_{k}}$ to the primal gap}.  
By construction of $d_{k}$ we have
\[
\inner{g_{k}}{d_{k}}
= \min\bigl\{\inner{g_{k}}{s_{k}-x_{k}},\,
\inner{g_{k}}{x_{k}-v_{k}}\bigr\}
\le \inner{g_{k}}{s_{k}-x_{k}}.
\tag{7}
\]
Moreover,
\[
\inner{g_{k}}{s_{k}-x_{k}}
= \inner{g_{k}}{s_{k}-x^{\star}}-\inner{g_{k}}{x_{k}-x^{\star}}
\le -\inner{g_{k}}{x_{k}-x^{\star}},
\tag{8}
\]
because $s_{k}$ solves the LMO and therefore
$\inner{g_{k}}{s_{k}}\le\inner{g_{k}}{x^{\star}}$.

Combining \eqref{7}–\eqref{8} yields
\[
-\inner{g_{k}}{d_{k}}\;\ge\;\inner{g_{k}}{x_{k}-x^{\star}} .
\tag{9}
\]

\noindent\textbf{[Step 5]} \emph{Using the geometric lower bound}.  
Lemma~\ref{lem:geom} gives
\[
\inner{g_{k}}{x_{k}-x^{\star}}
\;\ge\;
\frac{\mu\delta^{2}}{D^{2}}\,
\bigl(f(x_{k})-f(x^{\star})\bigr).
\tag{10}
\]

\noindent\textbf{[Step 6]} \emph{Insert \eqref{10} into \eqref{9} and then into \eqref{6}}.  
From \eqref{9} and \eqref{10},
\[
-\inner{g_{k}}{d_{k}}\;\ge\;
\frac{\mu\delta^{2}}{D^{2}}\,
\bigl(f(x_{k})-f(x^{\star})\bigr).
\tag{11}
\]
Squaring both sides and using the fact that the left‑hand side is non‑negative,
\[
\inner{g_{k}}{d_{k}}^{2}
\;\ge\;
\frac{\mu^{2}\delta^{4}}{D^{4}}\,
\bigl(f(x_{k})-f(x^{\star})\bigr)^{2}.
\tag{12}
\]

Plug \eqref{12} into \eqref{6}:
\[
f(x_{k+1})
\le f(x_{k})
-\frac{1}{2L D^{2}}\,
\frac{\mu^{2}\delta^{4}}{D^{4}}\,
\bigl(f(x_{k})-f(x^{\star})\bigr)^{2}.
\tag{13}
\]

\noindent\textbf{[Step 7]} \emph{From quadratic decrease to linear rate}.  
Define $h_{k}:=f(x_{k})-f(x^{\star})\ge0$.  Inequality \eqref{13} reads
\[
h_{k+1}\le h_{k}
-\underbrace{\frac{\mu^{2}\delta^{4}}{2L D^{6}}}_{\displaystyle\alpha}\,
h_{k}^{2}.
\tag{14}
\]
Since $h_{k}\le h_{0}$ for all $k$, we have $h_{k}^{2}\ge \frac{h_{k}}{h_{0}}\,h_{k}^{2}$,
which implies
\[
h_{k+1}\le h_{k}\Bigl(1-\alpha\frac{h_{k}}{h_{0}}\Bigr)
\le h_{k}\bigl(1-\rho\bigr),
\tag{15}
\]
where we set
\[
\rho:=\frac{\mu\delta^{2}}{L D^{2}}\in(0,1).
\tag{16}
\]
Iterating \eqref{15} yields
\[
h_{k}\le (1-\rho)^{k}h_{0},
\]
which is exactly the claim \eqref{1} of Theorem~\ref{thm:linear}. \hfill $\square$

\section*{Rate Derivation (With Constants)}
The constant $\rho$ in \eqref{16} can be expressed equivalently as
\[
\rho = \frac{\mu}{L}\,\frac{\delta^{2}}{D^{2}}
= \frac{1}{\kappa}\,\frac{\delta^{2}}{D^{2}},
\qquad
\kappa:=\frac{L}{\mu}
\text{ (condition number).}
\]
Thus the linear convergence factor improves when the polytope is ``well‑conditioned’’
(large $\delta$) or when $f$ is strongly convex (small $\kappa$).

\section*{Special Cases}
\begin{enumerate}
\item \textbf{Simplex.} For $\Cset$ being the probability simplex,
$\delta=1$ and $D=\sqrt{2}$, giving $\rho=\frac{2\mu}{L}$.
\item \textbf{Only convex (no strong convexity).} If $\mu=0$, the bound \eqref{14}
reduces to $h_{k+1}\le h_{k}$, and a standard curvature argument yields the classical
$O(1/k)$ rate:
\[
h_{k}\le \frac{2C_{f}}{k+2}.
\]
\end{enumerate}

\begin{thebibliography}{9}
\bibitem{lacoste2015linear}
S. Lacoste-Julien and M. Jaggi.
\newblock {Linear Convergence of Frank–Wolfe Variants}.
\newblock \emph{Adv. Neural Inf. Process. Syst.}, 28:33--41, 2015.
\end{thebibliography}

\end{document}